% =============================================================================
% Chapter 4: Implementation and Deployment  (14–18 pages target)
% =============================================================================
\chapter{Implementation and Deployment}
\label{chap:implementation}

\section*{Introduction}

This chapter details the concrete realization of the \eduspace{} platform.
It begins by describing the technological framework and development environment, covering both the hardware and software tools employed as well as the supporting configuration.
It then presents the DevOps and security practices adopted during development, including source code management, application security measures, and the design patterns leveraged throughout the codebase.
The chapter concludes with a presentation of the main user interfaces and their corresponding functionalities.

% =============================================================================
\section{Technological Framework and Development Environment}
\label{sec:tech-framework}

% ─────────────────────────────────────────────────────────────────────────────
\subsection{Hardware and Software Environment}
\label{sec:environment}

Development was carried out on personal machines running Linux (Ubuntu).
The following software tools and technologies were used across the project.

\subsubsection{Frontend Technologies}

\tabref{tab:frontend-stack} lists the frontend technology stack.

\begin{table}[H]
  \centering
  \caption{Frontend technology stack.}
  \label{tab:frontend-stack}
  \small
  \begin{tabularx}{\textwidth}{@{} l l X @{}}
    \toprule
    \textbf{Technology} & \textbf{Version} & \textbf{Purpose} \\
    \midrule
    React                & 19        & UI library for building the single-page application \\
    TypeScript           & 5         & Static type safety across the entire frontend codebase \\
    Vite                 & 6         & Build tool and development server with hot module replacement \\
    Tailwind CSS         & 4         & Utility-first CSS framework for responsive styling \\
    shadcn/ui + Radix UI & ---       & Accessible, composable component library \\
    React Router         & 7         & Client-side routing and navigation \\
    TanStack React Query & 5         & Server state management, caching, and data fetching \\
    TanStack React Table & 8         & Headless table library for the grades table \\
    FullCalendar         & 6         & Calendar component for assignment due date views \\
    Axios                & 1         & HTTP client for API communication \\
    Socket.IO Client     & 4         & WebSocket client for real-time group chat \\
    Lucide React         & ---       & Icon library \\
    \bottomrule
  \end{tabularx}
\end{table}

\subsubsection{Backend Technologies}

\tabref{tab:backend-stack} lists the backend technology stack, shared across all seven services and the API Gateway.

\begin{table}[H]
  \centering
  \caption{Backend technology stack.}
  \label{tab:backend-stack}
  \small
  \begin{tabularx}{\textwidth}{@{} l l X @{}}
    \toprule
    \textbf{Technology} & \textbf{Version} & \textbf{Purpose} \\
    \midrule
    Node.js          & 20 LTS    & JavaScript runtime for all backend services \\
    Express          & 5         & HTTP framework for REST API routing and middleware \\
    TypeScript       & 5         & Static type safety across all backend services \\
    Prisma ORM       & 6         & Type-safe database access, schema declaration, and migrations \\
    jsonwebtoken     & 9         & JWT creation and verification for authentication \\
    Multer           & 2         & Middleware for handling multipart file uploads \\
    Nodemailer       & 6         & Email sending for password reset and notification emails \\
    Morgan           & 1         & HTTP request logging middleware \\
    Socket.IO        & 4         & WebSocket server for real-time chat (Communication Service) \\
    amqplib          & 0.10      & RabbitMQ client for publishing and consuming events \\
    ioredis          & 5         & Redis client for caching, rate limiting, and Pub/Sub \\
    \bottomrule
  \end{tabularx}
\end{table}

\subsubsection{Infrastructure Technologies}

\tabref{tab:infra-stack} lists the infrastructure components, all orchestrated via Docker Compose.

\begin{table}[H]
  \centering
  \caption{Infrastructure technology stack.}
  \label{tab:infra-stack}
  \small
  \begin{tabularx}{\textwidth}{@{} l l X @{}}
    \toprule
    \textbf{Technology} & \textbf{Version} & \textbf{Purpose} \\
    \midrule
    PostgreSQL      & 16       & Relational database (five isolated databases) \\
    Redis           & 7        & Caching, rate limiting, presence tracking, Pub/Sub \\
    RabbitMQ        & 3.13     & Message broker for asynchronous inter-service events \\
    MinIO           & latest   & S3-compatible object storage for all binary files \\
    Elasticsearch   & 8        & Full-text search engine for post indexing and querying \\
    NGINX           & 1.25     & Reverse proxy unifying all services behind port 5000 \\
    pgAdmin         & 4        & Web-based PostgreSQL administration interface \\
    Docker          & 27       & Containerization of all infrastructure components \\
    Docker Compose  & 2        & Multi-container orchestration and configuration \\
    \bottomrule
  \end{tabularx}
\end{table}

% ─────────────────────────────────────────────────────────────────────────────
\subsection{Configuration and Support Tools}
\label{sec:tools}

\subsubsection{Integrated Development Environment}

Development was performed using \textbf{Visual Studio Code} with the following extensions: ESLint, Prettier, Prisma, Tailwind CSS IntelliSense.

\subsubsection{Version Control}

Source code was managed with \textbf{Git}~\cite{devops2024mslearn} and hosted on a private \textbf{GitHub} repository.
The team adopted a feature-branch workflow:
\begin{itemize}[nosep]
  \item The \texttt{master} branch contained the stable, working version of the project.
  \item Feature branches were created for each service or feature.
  \item Changes were integrated into \texttt{master} through pull requests, enabling code review before merging.
\end{itemize}

\subsubsection{Package Management}

The monorepo was managed using \textbf{pnpm} with \textbf{pnpm workspaces}.

\subsubsection{Database Management}

Database schemas were managed declaratively through \textbf{Prisma} schema files.


% =============================================================================
\section{DevOps Practices and Security}
\label{sec:devops}

% ─────────────────────────────────────────────────────────────────────────────
\subsection{Source Code Management and Git Workflow}
\label{sec:git-workflow}

The project followed a \textbf{feature-branch workflow} with the following conventions:

\begin{enumerate}[nosep]
  \item \textbf{Branch creation.}
        For each new service or feature, a dedicated branch was created from \texttt{master} using a descriptive naming convention (e.g., \texttt{feature/user-service}, \texttt{fix/avatar-upload-url}).

  \item \textbf{Development and commits.}
        Development proceeded on the feature branch with commit messages.

  \item \textbf{Pull request and review.}
        Upon completion, the developer opened a pull request on GitHub describing the changes.
        The other team member reviewed the code, suggested improvements, and approved the merge.

  \item \textbf{Merge and cleanup.}
        Approved pull requests were merged into \texttt{master} using a merge commit.
        The feature branch was deleted after merging to keep the repository clean.
\end{enumerate}


% Note: Section 4.2.2 "Pipeline CI/CD" is facultatif — omitted per instructions.
% Note: Section 4.2.3 "Tests automatisés" is facultatif — omitted per instructions.

% ─────────────────────────────────────────────────────────────────────────────
\subsection{Application Security}
\label{sec:security}

Security was a cross-cutting concern addressed at multiple layers of the architecture.

\subsubsection{Authentication and Token Management}

User authentication is based on \textbf{JSON Web Tokens} (JWT)~\cite{jwt2024rfc}.
Upon successful login, the User Service issues a signed JWT with a seven-day expiry containing the user's ID and email.
The token is stored client-side and sent in the \texttt{Authorization} header of every subsequent request.

The \textbf{API Gateway} intercepts all incoming requests, verifies the JWT signature and expiry, and extracts the authenticated user's identity.
If the token is invalid or expired, the request is rejected with a 401 status before it reaches any downstream service.
The verified user ID and email are forwarded as HTTP headers (\texttt{x-user-id}, \texttt{x-user-email}) to the downstream service, which trusts this information as pre-authenticated.

For WebSocket connections, the JWT is provided during the Socket.IO handshake.
The Communication Service validates the token before accepting the connection.

\subsubsection{Password Security}

User passwords are hashed using \textbf{bcrypt} before storage.

Password reset tokens are generated as cryptographically random strings, hashed (SHA-256) before storage, and expire after a configurable time window.

\subsubsection{Authorization and Access Control}

Authorization is enforced at the service level through role-based checks:

\begin{itemize}[nosep]
  \item \textbf{Classroom membership verification.}
        Before any operation on a classroom's content, the Content Service and Communication Service verify with the Class Service that the requesting user is an active member of the target classroom.
  \item \textbf{Role-based permissions.}
        Operations such as creating posts, managing chapters, grading submissions, and modifying classroom settings are restricted to users with the \texttt{ADMIN} role in the target classroom.
  \item \textbf{Creator-only operations.}
        Classroom deletion is restricted to the original creator, not all administrators.
  \item \textbf{Classroom type enforcement.}
        The API rejects grading-related requests (submission creation, grading, grade table access) for \emph{Friendly} classrooms, regardless of the user's role.
\end{itemize}

\subsubsection{File Access Control}

Files stored in MinIO are not publicly accessible.
Access is controlled through \textbf{presigned URLs} generated by the File Service on demand.
Each presigned URL is time-limited (expires after a configured duration), ensuring that file access requires a valid, non-expired link obtained through an authenticated API call.

\subsubsection{Internal Service Communication}

Certain API endpoints are designated as \textbf{internal} and are not exposed through the API Gateway's routing configuration.
These include endpoints for creating notifications in bulk, managing chat rooms, and indexing search content.
They are accessible only through direct service-to-service calls within the Docker network.

\figref{fig:security-layers} summarizes the security measures across the architecture layers.

\begin{figure}[H]
  \centering
  \begin{tikzpicture}[
    node distance=0.6cm,
    every node/.style={font=\small},
    layer/.style={draw, rounded corners, minimum width=13cm, minimum height=1cm, align=center},
    arrow/.style={-{Stealth[length=2.5mm]}, thick}
  ]
    \node[layer, fill=red!10]    (l1) {\textbf{NGINX} --- TLS termination, reverse proxy};
    \node[layer, fill=orange!10, below=of l1] (l2) {\textbf{API Gateway} --- JWT verification, rate limiting (Redis)};
    \node[layer, fill=yellow!10, below=of l2] (l3) {\textbf{Services} --- Role-based authorization, membership checks, type enforcement};
    \node[layer, fill=green!10,  below=of l3] (l4) {\textbf{Data Layer} --- Bcrypt passwords, hashed reset tokens, presigned file URLs};

    \draw[arrow] (l1) -- (l2);
    \draw[arrow] (l2) -- (l3);
    \draw[arrow] (l3) -- (l4);

  \end{tikzpicture}
  \caption{Security measures organized by architecture layer.}
  \label{fig:security-layers}
\end{figure}

% ─────────────────────────────────────────────────────────────────────────────
\subsection{Design Patterns}
\label{sec:design-patterns}

Several well-established design patterns were applied throughout the \eduspace{} codebase.

\begin{enumerate}[label=\textbf{DP-\arabic*.}]
  \item \textbf{Middleware Chain.}
        Express middleware functions are composed in a pipeline: authentication $\rightarrow$ validation $\rightarrow$ controller.
        Shared middleware (\texttt{authenticate}, \texttt{validate}, \texttt{errorHandler}) is centralized in the \texttt{@repo/shared} package.

  \item \textbf{Singleton.}
        Infrastructure clients (Prisma, Redis, RabbitMQ connections) are instantiated once and reused across the service lifetime, avoiding connection overhead.

  \item \textbf{Observer (Real-Time).}
        The Communication Service uses Socket.IO to implement the observer pattern: when a user sends a message, all connected clients in the same room are notified in real time via event broadcasting.
\end{enumerate}

% =============================================================================
\section{User Interfaces and Main Functionalities}
\label{sec:interfaces}

This section presents the main user interfaces of \eduspace{} as implemented in the React frontend.

\begin{figure}[H]
  \centering
  
\includegraphics[width=\textwidth]{screenshots/login-page.png}
  \caption{Login page with email/password form and branding panel.}
  \label{fig:ui-login}
\end{figure}

\begin{figure}[H]
  \centering
  \includegraphics[width=\textwidth]{screenshots/register-page.png}
  \caption{Registration page with form fields for creating a new account.}
  \label{fig:ui-register}
\end{figure}

\begin{figure}[H]
  \centering
  \includegraphics[width=\textwidth]{screenshots/dashboard.png}
  \caption{Dashboard showing classroom cards with type badges, member count, and quick actions.}
  \label{fig:ui-dashboard}
\end{figure}

\begin{figure}[H]
  \centering
  \includegraphics[width=0.55\textwidth]{screenshots/create-classroom.png}
  \caption{Create classroom dialog with name, type selection, and optional fields.}
  \label{fig:ui-create-class}
\end{figure}

\begin{figure}[H]
  \centering
  \includegraphics[width=0.55\textwidth]{screenshots/join-classroom.png}
  \caption{Join classroom dialog where users enter a 6-character class code.}
  \label{fig:ui-join-class}
\end{figure}

\begin{figure}[H]
  \centering
  \includegraphics[width=\textwidth]{screenshots/assignments-tab.png}
  \caption{Materials tab displaying posts organized by chapter in an accordion layout.}
  \label{fig:ui-materials}
\end{figure}

\begin{figure}[H]
  \centering
  \includegraphics[width=0.6\textwidth]{screenshots/create-post.png}
  \caption{Post creation form with type selector, chapter assignment, and file attachments.}
  \label{fig:ui-create-post}
\end{figure}

\begin{figure}[H]
  \centering
  \includegraphics[width=\textwidth]{screenshots/people-tab.png}
  \caption{People tab listing classroom members grouped by role (Teachers and Students).}
  \label{fig:ui-people}
\end{figure}

\begin{figure}[H]
  \centering
  \includegraphics[width=\textwidth]{screenshots/grades-tab.png}
  \caption{Grades tab with the grade table showing assignments, scores, and averages.}
  \label{fig:ui-grades}
\end{figure}

\begin{figure}[H]
  \centering
  \includegraphics[width=\textwidth]{screenshots/chat-tab.png}
  \caption{Real-time group chat interface with message history and online indicators.}
  \label{fig:ui-chat}
\end{figure}

\begin{figure}[H]
  \centering
  \includegraphics[width=\textwidth]{screenshots/account-settings.png}
  \caption{Account settings page for updating profile information.}
  \label{fig:ui-settings}
\end{figure}

\begin{figure}[H]
  \centering
  \includegraphics[width=\textwidth]{screenshots/security-settings.png}
  \caption{Security settings page for changing password.}
  \label{fig:ui-security}
\end{figure}

% =============================================================================
\section*{Conclusion}

This chapter detailed the implementation and deployment of \eduspace{}.
The technological framework was presented comprehensively, covering the frontend stack (React~19, TypeScript, Tailwind~CSS, TanStack libraries, FullCalendar, Socket.IO), the backend stack (Node.js, Express~5, Prisma, JWT, RabbitMQ, Redis), and the containerized infrastructure (PostgreSQL, Redis, RabbitMQ, MinIO, Elasticsearch, NGINX).
The development workflow, anchored by Git and GitHub with a feature-branch strategy, was described alongside the pnpm-based monorepo management and Prisma migrations.

Security was addressed as a layered concern---from NGINX and API Gateway-level protections through service-level authorization (role checks, membership verification, classroom type enforcement) to data-level safeguards (bcrypt hashing, presigned URLs, hashed reset tokens).